\section{Data}\label{sec:data}

This study targets the Northeast Massachusetts / Boston (NEMA) load zone of
ISO New England (ISO-NE), identified as load zone \texttt{.Z.NEMASSBOST} and
ISO-NE location id 4008. Three data sources are combined: historical and live
zonal load, weather observations and forecasts, and an operational day-ahead
demand benchmark.

\subsection{Load data}

The training target is the ISO-NE hourly wholesale-load series (RTLO) for NEMA,
covering 2017-03-01 through 2025-11-30, comprising 76{,}719 hourly observations
after cleaning. For live serving, load is drawn from the ISO-NE Web Services API
endpoint \texttt{realtimehourlydemand} (location 4008), which publishes with a
settlement lag of approximately four days; hourly demand from the U.S. Energy
Information Administration (EIA) is used as a fallback when ISO-NE data are
unavailable~\cite{eia930}. The two series differ slightly, as the
wholesale-load (WHLSECOST) history and the real-time demand series are not
identical products.

\subsection{Weather data}

Weather is obtained from Open-Meteo~\cite{openmeteo} for the Boston Logan station
(42.3656$^{\circ}$\,N, $-71.0096^{\circ}$\,W). The variables used are
temperature, relative humidity, wind speed, dew point, cloud cover, apparent
temperature, and visibility (in $^{\circ}$F, mph, and \%). Critically, a single
provider supplies both training history and live serving: the Open-Meteo archive
API (backed by ERA5 reanalysis) furnishes the historical record, while the
Open-Meteo forecast API supplies weather at serving time. Neither requires an API
key. This deliberate train/serve source unification removes a distribution
mismatch that otherwise degrades deployed day-ahead accuracy. The Open-Meteo
archive omits visibility, which is filled with a constant.

\subsection{Benchmark}

As an operational baseline we use ISO-NE's official day-ahead demand forecast for
NEMA, retrieved from the ISO-NE Web Services \texttt{dayaheadhourlydemand}
endpoint (location 4008)~\cite{isone}. This is the forecast that ISO-NE operators
rely on for unit commitment and reserve planning, and constitutes a meaningful,
hard bar for any candidate model.

\subsection{Temporal split, cleaning, and imputation}

To prevent any leakage of future information into model fitting, the data are
divided by a strict temporal split at \texttt{TRAIN\_CUTOFF}
$=$ 2024-12-31 23:00. All observations at or before the cutoff form the training
partition (68{,}704 rows) and all later observations form the held-out test
partition (8{,}015 rows). Because the model consumes fixed-length lookback
windows (Section~\ref{sec:method}), sequence construction yields 58{,}236
training, 10{,}276 validation, and 7{,}823 test windows, where the validation
set is the most recent 15\% of the training partition. Table~\ref{tab:data}
summarizes the partition sizes.

Cleaning and imputation are performed using statistics estimated on the training
partition only: missing values are filled with train-only medians, so that no
test-period information influences the imputation of any feature. Combined with
the temporal split, this ensures that every fitting decision depends solely on
information available at or before the cutoff.

\begin{table}[t]
\centering
\caption{NEMA load data partitions under the strict temporal split at
2024-12-31 23:00.}
\label{tab:data}
\begin{tabular}{lrr}
\toprule
Partition & Rows & Windows \\
\midrule
Train      & 68{,}704 & 58{,}236 \\
Validation & --       & 10{,}276 \\
Test       & 8{,}015  & 7{,}823 \\
\midrule
Total      & 76{,}719 & --       \\
\bottomrule
\end{tabular}
\end{table}
